% ============================================================
% VERSIÓN INDIVIDUAL — DANIEL DUARTE CORDERO
% Basada en la propuesta auditada del equipo + literatura asignada.
% ============================================================
\section{Metodología Propuesta}
\label{sec:metodologia}

\subsection{Formalización del problema}
\label{subsec:formalizacion}
El problema se formula como clasificación binaria probabilística sobre
datos tabulares. Sea $x_t\in\mathcal{X}$ el vector de características
construido exclusivamente con información observable hasta el instante
$t$, y sea $y\in\{0,1\}$ el resultado final desde la perspectiva del
jugador al que corresponde la fila. Se aprende una función
$f_\theta:\mathcal{X}\rightarrow[0,1]$ que aproxime la probabilidad
condicional de victoria.

\begin{equation}
\label{eq:objetivo}
  f_\theta(x_t) \approx P(y=1\mid x_t),
  \qquad
  \theta^{*}=\arg\min_{\theta}\sum_{i\in\mathcal{D}_{\mathrm{train}}}
  \ell\!\left(y_i,f_\theta(x_i)\right),
\end{equation}

donde $\ell$ es una función de pérdida probabilística apropiada para el
modelo y toda selección de hiperparámetros se realiza sin utilizar el
conjunto de prueba.

\noindent\textbf{Definición del target.} La etiqueta se deriva de
\texttt{rewards} a nivel de episodio: $y=1$ si gana el jugador desde
cuya perspectiva se construye el snapshot y $y=0$ si pierde. La
auditoría resolvió el resultado en 99.86\% de los episodios. Los 33
empates se excluyen inicialmente del target binario. Los 73 finales por
\texttt{TIMEOUT}, \texttt{ERROR} o \texttt{INVALID} se excluyen del
entrenamiento principal y se reservan para análisis de sensibilidad,
porque el desenlace no procede de la dinámica normal del juego.

\noindent\textbf{Unidad de observación.} Una fila corresponde a un
snapshot fresco del asiento activo, proyectado a una perspectiva común
\emph{self}/\emph{opp}. La observación congelada del asiento inactivo no
se duplica como fila adicional.

\noindent\textbf{Unidad independiente.} La unidad estadísticamente
independiente es el episodio. Los 52,085 episodios auditados generan
más de siete millones de filas, pero ningún episodio puede aparecer en
más de una partición. Esta distinción gobierna splits, validación e
intervalos de confianza.

\subsection{Pipeline previsto}
\label{subsec:pipeline}
El pipeline ingiere los replays JSON, valida la firma de esquema y
\texttt{schema\_version}, selecciona los estados frescos del asiento
activo y normaliza la perspectiva mediante \texttt{yourIndex}. Luego
construye características de progreso, recursos, tablero, energía,
estados especiales y diferenciales entre ambos jugadores. Las columnas
\texttt{episode\_id}, fecha y turno se conservan para agrupación,
estratificación y análisis, pero no entran al vector predictivo. La
tabla materializada se divide por episodio antes de ajustar cualquier
transformación. Sobre train se ajustan preprocesamiento y modelos; la
validación selecciona hiperparámetros y el test queda intacto hasta la
evaluación final. En Etapa 3, el mismo extractor convertirá un estado
actual en features y entregará $P(\mathrm{win}\mid\mathrm{state})$ al
agente como evaluador de posiciones.

\begin{figure}[!t]
  \centering
  \safeimage[width=\columnwidth]{pipeline.png}
  \caption{Pipeline propuesto: replay JSON $\rightarrow$ validación de
  esquema $\rightarrow$ snapshots frescos normalizados $\rightarrow$
  tabla de features $\rightarrow$ particiones por episodio $\rightarrow$
  modelos y evaluación $\rightarrow$ probabilidad consumida por el agente.}
  \label{fig:pipeline}
\end{figure}

\subsection{Protocolo experimental}
\label{subsec:protocolo}

\subsubsection{Particiones}
El protocolo principal será 70/15/15 para train, validación y test,
\textbf{agrupando por \texttt{episode\_id}}. Se descarta el split por
snapshot porque estados de una misma partida son fuertemente
correlacionados y su presencia a ambos lados de la partición produciría
fuga. Como protocolo secundario se usará una partición temporal:
train con 2026-06-17, 06-18, 06-23 y 06-29; validación con 07-05 y
07-11; test con 07-25 y 08-04. El día 06-16 se mantiene fuera por ser
un arranque parcial y 07-17 se reserva por su concentración extrema de
mazos/agentes, sujeto a revisión posterior al EDA. Un tercer análisis de
robustez reservará mazos o pares de mazos completos para prueba. El
objetivo de estos protocolos separados es distinguir nueva partida,
futuro temporal y composición no vista, una preocupación coherente con
la evaluación entre parches y equipos observada en literatura
\cite{chitayat_meta_2023,angliss_vgcbench_2026}.

\subsubsection{Estrategia de validación}
La selección de hiperparámetros se realizará exclusivamente dentro de
train mediante validación cruzada agrupada por episodio
(\texttt{GroupKFold}) o una variante equivalente que preserve grupos.
La partición de validación externa se usa para comparar configuraciones
y fijar decisiones antes de la prueba final. El test principal se toca
una sola vez para la evaluación reportada. En el protocolo temporal se
respeta el orden cronológico y no se reutilizan datos futuros para
ajustar transformaciones.

\subsubsection{Métricas y su justificación}
\begin{table}[!t]
\centering
\caption{Métricas de evaluación y su justificación.}
\label{tab:metricas}
\footnotesize
\begin{tabularx}{\columnwidth}{@{}P{2.2cm} Y@{}}
\toprule
\textbf{Métrica} & \textbf{Por qué esta y qué alternativa desplaza}\\
\midrule
ROC-AUC & Métrica principal de discriminación, independiente de un
umbral y directamente comparable con el antecedente de Hearthstone
\cite{janusz_hearthstone_2017}.\\
Log loss & Evalúa la calidad de la probabilidad y penaliza con fuerza
predicciones seguras y equivocadas; complementa lo que AUC no mide.\\
Brier score + calibración & Mide error cuadrático de probabilidad y se
acompaña con curvas de fiabilidad; es necesario porque una corrección de
clases puede preservar discriminación y deteriorar calibración
\cite{vandenGoorbergh_imbalance_2022}.\\
\bottomrule
\end{tabularx}
\end{table}

ROC-AUC es principal porque conecta con el único antecedente cercano y
responde si el modelo ordena estados ganadores por encima de perdedores.
No obstante, una salida destinada a interpretarse como probabilidad no
puede evaluarse solo por ranking o accuracy. Van den Goorbergh et al.
muestran que intervenciones de balance pueden dejar el AUROC
prácticamente sin mejora y, simultáneamente, distorsionar las
probabilidades \cite{vandenGoorbergh_imbalance_2022}. Por eso se
reportan log loss, Brier score y curvas de calibración, globalmente y
por ventana de turno. Los intervalos de confianza se estimarán por
bootstrap sobre episodios, no sobre snapshots.

\subsection{Plan de baselines}
\label{subsec:planbaselines}
Cada nivel añade complejidad solo para responder una pregunta que el
anterior no puede resolver.

\begin{table}[!t]
\centering
\caption{Escalera de baselines.}
\label{tab:baselines}
\footnotesize
\begin{tabularx}{\columnwidth}{@{}c P{2.4cm} Y@{}}
\toprule
\textbf{Nivel} & \textbf{Modelo} & \textbf{Pregunta que responde}\\
\midrule
B0 & \texttt{DummyClassifier} (prior / clase mayoritaria) & ¿Cuál es el piso de rendimiento?\\
B1 & Heurística simple de recursos o premios, fijada tras EDA & ¿Hace falta ML o basta una regla de dominio?\\
B2 & Regresión logística regularizada & ¿La relación es esencialmente lineal y ofrece una referencia probabilística interpretable?\\
B3 & Árbol de decisión / Random Forest & ¿Hay no linealidades e interacciones que B2 no captura?\\
B4 & Gradient Boosting (familia por fijar) & ¿Cuánto rendimiento adicional compra la complejidad del ensamble?\\
\bottomrule
\end{tabularx}
\end{table}

\noindent\textbf{Tercera referencia.} El rango 0.7846--0.8019 AUC de
Hearthstone se arrastra como referencia externa parcialmente comparable
junto con los baselines internos \cite{janusz_hearthstone_2017}.

\subsection{Contrato de reproducibilidad}
\label{subsec:reproducibilidad}
La auditoría ya utiliza la semilla fija \texttt{20260821}. La etapa de
modelado mantendrá las semillas y configuraciones en archivos
versionados, fijará dependencias, registrará cada ejecución y su commit,
y guardará tablas comparables de parámetros y métricas. El repositorio
documentará hardware y el comando exacto de reproducción de cada
resultado una vez el pipeline quede implementado. \textbf{\% TODO
implementación/integración:} insertar aquí los nombres finales de los
archivos de configuración y comandos reproducibles; no se inventan en
esta etapa.

\subsection{Riesgos y mitigación}
\label{subsec:riesgos}
\begin{table}[!t]
\centering
\caption{Riesgos identificados y su mitigación.}
\label{tab:riesgos}
\footnotesize
\begin{tabularx}{\columnwidth}{@{}P{2.0cm} c Y@{}}
\toprule
\textbf{Riesgo} & \textbf{Impacto} & \textbf{Mitigación / disparador}\\
\midrule
Fuga entre snapshots del mismo episodio & A & Split y CV obligatoriamente agrupados por \texttt{episode\_id}; prueba de intersección vacía.\\
Uso de información futura/privada & A & Lista negra de \texttt{rewards}, \texttt{visualize}, metadata y campos posteriores al snapshot; test automático de columnas.\\
Cambio temporal de población/versión & A & Reportar holdout agrupado y temporal por separado; tratar la caída como resultado y no ocultarla.\\
Memorización de mazos frecuentes & M/A & Holdout por mazo/matchup y ablación de contexto de mazo; interpretar ganancias de ID con cautela.\\
\bottomrule
\end{tabularx}
\end{table}
